\documentclass[11pt]{article}

\usepackage[T1]{fontenc}
\usepackage[utf8]{inputenc}
\usepackage{amsmath,amsthm,amssymb}
\usepackage[margin=1.1in]{geometry}
\usepackage[colorlinks=true,linkcolor=blue,citecolor=blue,urlcolor=blue]{hyperref}

\newtheorem{theorem}{Theorem}[section]
\newtheorem{proposition}[theorem]{Proposition}
\newtheorem{lemma}[theorem]{Lemma}
\newtheorem{corollary}[theorem]{Corollary}
\theoremstyle{definition}
\newtheorem{definition}[theorem]{Definition}
\theoremstyle{remark}
\newtheorem{remark}[theorem]{Remark}

\newcommand{\C}{\mathbb{C}}
\newcommand{\R}{\mathbb{R}}
\newcommand{\Q}{\mathbb{Q}}
\newcommand{\Z}{\mathbb{Z}}
\newcommand{\im}{\operatorname{im}}
\newcommand{\Frac}{\operatorname{Frac}}
\newcommand{\val}{\operatorname{val}}
\DeclareMathOperator{\antidiag}{antidiag}

\emergencystretch=1.5em

\title{A zero-dimensional calibration of field redefinitions for AQFT
and pAQFT: the Alp\"oge--Mathew counterexample and uniqueness of its
defect}

\author{Christoph Mayer\thanks{Independent researcher.
Email: \texttt{hcab14@gmail.com}.}}

\date{July 25, 2026}

\begin{document}

\maketitle

\begin{abstract}
The Alp\"oge--Mathew map $F:\C^3\to\C^3$ (July 2026) has constant
Jacobian determinant $-2$ but is not injective, disproving the Jacobian
Conjecture for $n\ge 3$. Read as a zero-dimensional QFT ---
$F(\varphi)=J$ a classical field equation, the Bass--Connell--Wright
tree expansion its perturbation theory --- it is an exactly solvable
toy model in which loops are trivial, the tree series \emph{converges},
and the theory nevertheless fails globally: two of three solution
sheets sit at infinity over the perturbative vacuum and are invisible
to every order of perturbation theory. Our primary interest is what
this calibrates for rigorous QFT, especially algebraic QFT (AQFT) and
perturbative algebraic QFT (pAQFT). Exact results, verified in a public
repository: (i)~the pullback $F^*$ is a monomorphism but not an
automorphism of the polynomial observable algebra --- $\{1,x,x^2\}$ is
a basis over $\Frac(\im F^*)$, and $x$ is not integral over $\im F^*$
(escape-curve certificate $\Leftrightarrow$ non-properness), the
cleanest 0D vindication of the Haag--Kastler preference for algebras
over field coordinates; (ii)~this calibrates the formal deferral in
pAQFT (Fredenhagen--Rejzner): ``upgrade formal to convergent'' is
\emph{not} the missing step to a non-perturbative construction ---
properness of the classical dynamics is; (iii)~an exact reduced Keller
identity for the weight family $(1,-1,-m)$; (iv)--(v)~complete
$v$-linear classifications in the sibling systems $(1,-1,-3)$ and
$(2,-1,-3)$, culminating in a uniqueness theorem: for every $m\ge 3$
the $v$-linear class of $(1,-1,-m)$ contains only tame automorphisms,
so the defect is rigidly attached to the $m=2$ numerology. We state
honestly what the model does \emph{not} imply (nothing about UV
renormalization, Borel summability, or $D=4$ existence). AI-assisted
exploratory research, not peer reviewed; every claim is labelled exact,
box-checked, or open, and is reproducible from a cited script.
\end{abstract}

\tableofcontents

%=====================================================================
\section{Introduction}
%=====================================================================

\subsection{Background}

The Jacobian Conjecture (Keller, 1939 \cite{Kel39}) asserted that every
polynomial map $F:\C^n\to\C^n$ with constant nonzero Jacobian determinant
(a \emph{Keller map}) is a polynomial automorphism. Despite the
degree-reduction theory of Bass--Connell--Wright \cite{BCW82}, the formal
tree-inversion formulas of Wright \cite{Wri87}, and decades of partial
results collected in van den Essen's monograph \cite{vdE00}, the
conjecture remained open in every dimension $n\ge 2$ until July 19, 2026,
when Levent Alp\"oge announced (with Akhil Mathew; the map itself was
produced by a large language model in response to Mathew's question) an
explicit counterexample in dimension $3$ \cite{AM26}: the map
\begin{equation}\label{eq:AM}
F(x,y,z) \;=\; \bigl(\,(1+xy)^3 z + y^2(1+xy)(4+3xy),\;\;
 y + 3x(1+xy)^2 z + 3xy^2(4+3xy),\;\;
 2x - 3x^2y - x^3 z\,\bigr)
\end{equation}
satisfies $\det DF \equiv -2$ yet identifies the three points
$(0,0,-\tfrac14)$, $(1,-\tfrac32,\tfrac{13}{2})$,
$(-1,\tfrac32,\tfrac{13}{2})$, all mapping to $(-\tfrac14,0,0)$. Hence
the conjecture is false for all $n\ge 3$ (adjoin identity coordinates);
the case $n=2$ remains open. The example was verified by the community
and formalized in Lean~4; all claims about it in this paper are
independently re-verified in the repository accompanying this paper
(\texttt{scripts/verify\_counterexample.py}).

The Jacobian Conjecture has a well-known reformulation in
zero-dimensional perturbative quantum field theory, made precise by
Abdesselam \cite{Abd03} on the basis of the tree formulas of
\cite{BCW82,Wri87} and the constructive combinatorics of
Abdesselam--Rivasseau \cite{AR95}: $F(\varphi)=J$ is the classical field
equation of a three-component scalar model in $D=0$ spacetime dimensions
with external source $J$; the linearization $DF(0)$ is the inverse
propagator; the nonlinear monomials of $F$ are interaction vertices; and
the perturbative inversion $\varphi(J)$ is exactly the sum over rooted
tree Feynman graphs. Constant Jacobian determinant means all loop
corrections are field-independent constants: the theory looks free at
loop level. A genuine counterexample therefore provides, for the first
time, an exactly solvable toy model in which perturbation theory is well
defined at every order, loop corrections are trivial, and the theory
still fails globally.

\emph{Why this matters for AQFT and pAQFT.} The motivation of the present
paper is not primarily to enlarge the catalogue of Jacobian
counterexamples, but to use this toy model as a \emph{calibration device}
for rigorous QFT --- especially the algebraic tradition of
Haag--Kastler \cite{HK64} and the perturbative algebraic QFT (pAQFT)
framework of Fredenhagen--Rejzner \cite{FR12,FR16,Rej16}. Polynomial
field redefinitions with $\det DF=\mathrm{const}\neq 0$ are used
throughout renormalization theory; invertibility is usually established
at the formal-power-series level, where it is automatic. The
counterexample shows that, already with three field components, a
polynomial redefinition with constant Jacobian need not be globally
injective: the formal inverse exists, \emph{converges} near the vacuum,
and is still only a local section of a 3-sheeted covering. The exact
structure of the observable algebra $\im F^*$ (Section~\ref{sec:obs})
turns the Haag--Kastler preference for algebras over field coordinates
into a theorem-sized 0D fact; the pAQFT deferral of convergence and
globality is calibrated rather than contradicted
(Section~\ref{sec:qftread}). The sibling searches and the uniqueness
theorem then ask how rigid this defect is: is it a generic feature of
constant-Jacobian dynamics, or a delicate coincidence of the $m=2$
weight arithmetic?

\subsection{Contents and main results}

This paper condenses, into a self-contained account, the exact-structure
results obtained in the public repository
\begin{center}
\url{https://github.com/hcab14/JacobianConjectureQFT}
\end{center}
about the Alp\"oge--Mathew map and its two nearest
$\C^*$-equivariance siblings, with the QFT/AQFT/pAQFT reading as the
primary interpretive frame. The main results are:

\begin{itemize}
\item \textbf{Section \ref{sec:qft}} sets up the $D=0$ QFT dictionary
 for the map \eqref{eq:AM}: the first-order (conjugate-field) action,
 the $\C^*$-grading with source weights $(1,-1,-2)$, the escape locus
 $\{p=0\}$ (the non-properness set), and the reduction of $\det DF$ to
 the two-dimensional invariant plane.
\item \textbf{Section \ref{sec:obs}} gives the exact structure of the
 observable algebra $\im F^*$: $\C(x,y,z)$ is free with basis
 $\{1,x,x^2\}$ over the fraction field of $\im F^*$, every polynomial
 observable has a unique computable normal form
 $c_0(F)+c_1(F)x+c_2(F)x^2$, and $x$ is \emph{not integral} over
 $\im F^*$ --- an exact escape-curve certificate showing that no
 finite-module refinement is possible (Theorems \ref{thm:basis},
 \ref{thm:nonintegral}).
\item \textbf{Section \ref{sec:redw}} states the exact reduced Keller
 identity for the whole weight family $(1,-1,-m)$:
 $\det DF = \det M$ and $R^m \det M = J_2(PR^m, QR)$, reducing the 3D
 Keller problem to two variables (Proposition \ref{prop:redw}).
\item \textbf{Section \ref{sec:213}} treats the weight system
 $(2,-1,-3)$, whose invariant ring is an $A_1$ quadric cone: a
 chart-based reduced Keller identity, a complete classification of the
 degree-$1$ Keller box (two families of tame automorphisms, Theorem
 \ref{thm:213class}), provable emptiness of the $\Z_2$ and $\Z_3$
 stabilizer-jump mechanisms in two boxes, and an honest account of the
 computational wall met by the degree-$2$ box.
\item \textbf{Section \ref{sec:113}} classifies the $v$-linear class of
 $(1,-1,-3)$ for all $w$-degrees by a Wronskian stratification
 (Theorem \ref{thm:113main}): every Keller map is a tame automorphism;
 the $3{:}1$ mechanism is empty; the Alp\"oge--Mathew mechanism is
 obstructed at $m=3$ by an exact numerological incompatibility
 (Section \ref{sec:numerology}). One substratum ($D_3$ with
 non-squarefree data) is a Gr\"obner-box gap (Section \ref{sec:gap}).
\item \textbf{Section \ref{sec:11m}} is the headline theorem: the
 $m=3$ classification extends uniformly to \emph{every} $m\ge 3$
 (Theorem \ref{thm:11m}). For $m\ge 4$ the D-strata are empty with no
 box at all; Alp\"oge--Mathew survives at $m=2$ exactly where the
 killing factor $u^{m-2}$ trivializes
 (\texttt{scripts/search\_11m.py}, \texttt{docs/SEARCH\_11M.md}).
\item \textbf{Section \ref{sec:qftread}} is the interpretive core: what
 the model calibrates for AQFT (Haag--Kastler), pAQFT
 (Fredenhagen--Rejzner), and non-perturbative algebraic constructions
 (Buchholz--Fredenhagen), together with an explicit ledger of what it
 does \emph{not} imply; the uniqueness theorem is read as rigidity of
 the non-properness defect.
\item \textbf{Sections \ref{sec:outlook} and \ref{sec:repro}} give
 QFT-oriented next probes and the reproducibility table.
\end{itemize}

\subsection{Provenance and verification}\label{sec:provenance}

This paper reports \textbf{AI-assisted exploratory research}: it was
produced in interactive sessions with AI coding agents (large language
model agents wrote and ran the verification scripts), and it has
\textbf{not been peer reviewed}. The counterexample itself is due to
Alp\"oge and Mathew \cite{AM26}; this work only studies it. The claims
should be read with the following calibration, which the text states
case by case:

\begin{itemize}
\item \emph{Exact / symbolic results} (the bulk of the paper): every
 displayed identity and every emptiness claim stated as proved is
 verified by exact symbolic computation --- sympy \cite{sympy} over
 $\Q$, and Gr\"obner-basis certificates computed with msolve
 \cite{msolve} and Singular \cite{Singular} --- in the cited scripts of
 the public repository. ``Verified'' means the script passes, within
 the stated scope; the scope lines (often box- or chart-limited)
 matter and are always given.
\item \emph{Numerical results} (e.g.\ the $S_3$ monodromy of the three
 sheets, rigidity continuation): labelled as evidence, not proof.
\item \emph{Interpretation} (the QFT and anomaly readings): clearly
 flagged as such.
\end{itemize}

Every theorem and proposition below cites the repository script that
asserts it (for the main theorem this is
\texttt{scripts/search\_113.py}); Section \ref{sec:repro} lists the
commands and runtimes needed to reproduce all of them. Corrections and
counter-arguments are welcome.

\subsection{Notation}\label{sec:notation}

Throughout, $F=(F_1,F_2,F_3):\C^3\to\C^3$ is a polynomial map in source
coordinates $\varphi=(x,y,z)$ with target (source-of-the-QFT)
coordinates $J=(a,b,c)$. $DF$ is the Jacobian matrix and a \emph{Keller
map} has $\det DF \equiv \kappa \in \C^\times$. For a weight system
$(1,-1,-m)$ we write
\[
w = xy, \qquad v = x^m z
\]
for the generators of the invariant ring, and
\[
J_2(A,B) \;=\; \partial_w A\,\partial_v B - \partial_v A\,\partial_w B
\]
for the Jacobian of two functions of $(w,v)$. Primes denote $d/dw$. For
the weight system $(2,-1,-3)$ the invariants are $u_1=xy^2$,
$u_2=x^2yz$, $u_3=x^3z^2$ with chart coordinates $(t,p,q)$ as in Section
\ref{sec:213}. The polynomials
\[
p \;=\; 27a^2c^2 - 18abc + b^3c - b^2 + 16a, \qquad
D_0 \;=\; 27ac^2 - 9bc + 8
\]
are the escape (non-properness) polynomial and the fiber-parametrization
discriminant factor of the Alp\"oge--Mathew map (the symbol $p$ is also
used for the chart coordinate $u_1$ in Section \ref{sec:213} and for
lowercase coefficient functions $p_0,p_1$ in Section \ref{sec:113}; the
context always disambiguates).

%=====================================================================
\section{The Alp\"oge--Mathew map as a zero-dimensional QFT}
\label{sec:qft}
%=====================================================================

\subsection{The counterexample and the dictionary}

Direct computation (\texttt{scripts/verify\_counterexample.py}) confirms
$\det DF \equiv -2$ for the map \eqref{eq:AM} and the triple fiber over
$(-\tfrac14,0,0)$. Note that all three colliding points are real: the
restriction $F:\R^3\to\R^3$ is a real polynomial local diffeomorphism
with \emph{constant} Jacobian that fails to be injective.

One would like to write $F(\varphi)=J$ as the stationarity condition of
an action $S(\varphi)$, but that is impossible: $F$ is a gradient iff
$DF$ is symmetric, and here already
$\partial F_1/\partial z|_0 = 1 \ne 2 = \partial F_3/\partial x|_0$.
The standard remedy \cite{Abd03} is a conjugate (auxiliary) field
$\bar\varphi=(\bar x,\bar y,\bar z)$ and the first-order action
\[
S(\bar\varphi,\varphi) = \bar\varphi\cdot\bigl(F(\varphi)-J\bigr),
\qquad
Z(J) = \int d^3\bar\varphi\, d^3\varphi\; e^{-S(\bar\varphi,\varphi)},
\]
where the $\bar\varphi$-integral localizes $Z(J)$ on
$\delta^3(F(\varphi)-J)$. The ghost sector that would produce
$\det DF(\varphi)$ is trivial here, contributing only the constant
$-2$. Splitting $S = S_{\rm free} + S_{\rm int}$ with
$S_{\rm free} = \bar\varphi\cdot L\varphi - \bar\varphi\cdot J$,
$L = DF(0)$, the interaction consists of thirteen vertices, each with
exactly one $\bar\varphi$-leg, from cubic ($\bar x\, y^2$, $\bar y\, xz$,
$\bar z\, x^2y$) to octic ($\bar x\, x^3y^3z$). The one-point function
$\langle\varphi\rangle = G(J)$ is the formal inverse of $F$: Wick
contraction generates precisely the Bass--Connell--Wright--Wright sum
over rooted trees \cite{BCW82,Wri87}, with no loop corrections beyond
the constant $\det DF = -2$.

Two facts about this expansion, established in the repository and used
below only as context, are worth recording. First, the tree series
\emph{converges}, with finite nonzero radius; it is the Taylor series of
one branch of an explicit degree-3 algebraic function: Gr\"obner
elimination (\texttt{scripts/branch\_locus.py}) shows that the
$x$-coordinate of any preimage of $(a,b,c)$ satisfies the eliminant
cubic
\begin{equation}\label{eq:cubic}
p\,x^3 + q\,x + r = 0,\qquad
q = 4-3bc, \qquad r = -2c,
\end{equation}
with $p$ as in Section \ref{sec:notation}. Second, what fails is
\emph{globality}, not summability: on $\{p=0\}$ the cubic drops degree
and a solution sheet escapes to infinity in field space. The set
$\{p=0\}$ is exactly the non-properness set of $F$ in the sense of
Jelonek \cite{Jel93}, the perturbative vacuum $J=0$ lies \emph{on} it,
and the two non-perturbative sheets over the vacuum sit at infinity,
invisible to every order of perturbation theory. Off $\{p=0\}$ the map
is a genuine 3-sheeted covering (being \'etale, sheets never merge at
finite points); its geometric monodromy is the full $S_3$
(\emph{numerical} evidence, two independent loops;
\texttt{scripts/monodromy.py}).

\subsection{\texorpdfstring{$\C^*$}{C*}-equivariance and reduction to
the invariant plane}
\label{sec:equivariance}

The model is equivariant under the $\C^*$-action
$\lambda.(x,y,z) = (\lambda x, \lambda^{-1}y, \lambda^{-2}z)$: the
components of $F$ scale with weights $(-2,-1,1)$, so the sources
$(a,b,c)$ carry weights $(-2,-1,1)$ and every term of the action has
weight $0$. The invariant ring of the source action is the \emph{free}
polynomial ring $\C[w,v]$ with $w=xy$, $v=x^2z$, and the
Alp\"oge--Mathew map is exactly of the equivariant normal form
\[
F = \Bigl(\frac{P(w,v)}{x^{2}},\ \frac{Q(w,v)}{x},\ x\,R(w,v)\Bigr),
\]
with
\[
P_{\rm AM} = (1+w)^3 v + w^2(1+w)(4+3w),\quad
Q_{\rm AM} = w + 3(1+w)^2 v + 3w^2(4+3w),\quad
R_{\rm AM} = 2 - 3w - v .
\]
For such maps $\det DF$ descends to a polynomial in $(w,v)$ alone, so
the 3D Keller condition is a 2D problem --- this is the special case
$m=2$ of Proposition \ref{prop:redw} below, and it is the engine of the
classification results in Sections \ref{sec:213} and \ref{sec:113}. The
counterexample mechanism is a \emph{stabilizer jump}: the target
$a$-axis (weight $-2$) has stabilizer $\Z_2$, so a free source orbit
mapping into it (a common zero of $(Q,R)$ with $P\ne0$) is $2{:}1$
--- realized for \eqref{eq:AM} along the explicit curve
$\varphi(s) = (\pm s, \mp\tfrac{3}{2s}, \tfrac{13}{2s^2})$ with
$F(\varphi(s)) = (-\tfrac{1}{4s^2},0,0)$.

%=====================================================================
\section{The exact structure of the observable algebra
\texorpdfstring{$\im F^*$}{im F*}}
\label{sec:obs}
%=====================================================================

The pullback $F^*:\C[a,b,c]\to\C[x,y,z]$, $F^*\mathcal{O} =
\mathcal{O}\circ F$, is injective ($F$ is dominant) and not surjective
($F$ is $3$-to-$1$): the polynomial observables of the ``redefined
field'' $F(\varphi)$ form a proper subalgebra. This section makes the
cokernel completely explicit. All statements are asserted exactly in
\texttt{scripts/missing\_observables.py} ($\sim$15 s).

\begin{theorem}[Field-level structure of $\im F^*$; verified]
\label{thm:basis}
Let $K = \C(a,b,c) \hookrightarrow L = \C(x,y,z)$ via $F^*$. Then:
\begin{enumerate}
\item $L = K(x)$: adjoining the single coordinate $x$ gives everything;
 $y$ and $z$ are explicit $K$-rational expressions in $x$ with
 denominators involving only $D_0(F)$ (the fiber parametrization of the
 repository, read backwards).
\item $[L:K]=3$ with basis $\{1,x,x^2\}$: $x$ satisfies the eliminant
 cubic \eqref{eq:cubic} with coefficients evaluated at $F$, and the
 cubic is irreducible over $K$.
\item \emph{Normal form:} every polynomial observable
 $\mathcal{O}(x,y,z)$ is uniquely
 \[
 \mathcal{O} \;=\; c_0(F) + c_1(F)\,x + c_2(F)\,x^2,
 \qquad c_i \in \C(a,b,c),
 \]
 computable by one polynomial division (reduction mod the eliminant);
 the denominators of the $c_i$ are powers of $D_0$ times powers of the
 escape polynomial $p$.
\item \emph{Membership criterion:} $\mathcal{O}$ is a function of the
 redefined field alone iff $c_1 = c_2 = 0$. The missing observables
 are, at field level, the free module generated by $\{x, x^2\}$ over
 $\im F^*$ --- precisely the two sheet-separating degrees of freedom of
 the degree-3 covering.
\end{enumerate}
\end{theorem}

Worked examples (asserted): $F_1 \mapsto (a,0,0)$;
$F_3^2 - 5F_2 \mapsto (c^2-5b,\,0,\,0)$; and
\[
y \;=\; \frac{81abc^2 - 72ac - 15b^2c + 16b}{2D_0}
 \;+\; \frac{3p}{D_0}\,x
 \;-\; \frac{3(9ac-b)\,p}{2D_0}\,x^2 ,
\]
with a similar expansion for $z$ whose separator coefficients also carry
the factor $p$. The escape polynomial appears \emph{verbatim} in the
separator coefficients: the amount by which a coordinate fails to be an
observable of $F(\varphi)$ is proportional to the non-properness
polynomial.

\begin{theorem}[No finite-module version; exact escape certificate;
verified]\label{thm:nonintegral}
$x$ is not integral over $\im F^*$. Along the curve
\[
\varphi(T) \;=\; \Bigl(T,\; y_0,\; \frac{2T - 3T^2y_0 - c_3}{T^3}\Bigr),
\]
the source escapes ($x = T \to \infty$) while $F(\varphi(T))$ converges
to the finite target
$\bigl(y_0^2(1-c_3y_0),\; y_0(4-3c_3y_0),\; c_3\bigr)$, which lies
exactly on the wall $\{p=0\}$ (this is a rational parametrization of the
Jelonek set by $(y_0,c_3)$). An integral element must stay bounded when
its coefficients do, so no monic relation for $x$ over $\C[a,b,c]$
exists; in particular $\C[x,y,z]$ is not a finite module over
$\im F^*$.
\end{theorem}

\begin{remark}
This is the ring-theoretic face of non-properness: for these dominant
Keller maps, finite module $\Leftrightarrow$ proper map, and a proper
Keller map would be injective. Every counterexample to the Jacobian
Conjecture must therefore exhibit this infinite-module pathology; the
normal form of Theorem \ref{thm:basis} is the best possible structure
statement.
\end{remark}

Membership can be decided cheaply: field-level membership by the normal
form (one division); strict polynomial membership by the standard
tag-variable Gr\"obner test; and for the coordinates themselves a single
exact fiber computation suffices --- the fiber over $F(1,2,3)$ consists
of three points whose $x$-, $y$- and $z$-coordinates are pairwise
distinct (verified), and any element of $\im F^*$ is constant on fibers.

\emph{QFT / AQFT reading.} The observable algebra of the redefined
field has corank $2$ as a module inside the full observable algebra,
with explicit generators $x, x^2$ for the quotient: the operators that
distinguish the three classical vacua over a source $J$. They are
exactly the observables that perturbation theory around one vacuum
constructs as algebraic functions of $J$ (branches of the cubic) and
that no polynomial function of $J$ represents; they cease to exist as
finite quantities exactly where $p = 0$. ``Which observables are
missing'' and ``where does the theory degenerate'' have the same
answer. In the language of Section~\ref{sec:aqft}: $F^*$ is a
monomorphism but not an automorphism of the observable algebra --- the
0D fact that underwrites the Haag--Kastler preference for algebras
over field coordinates.

%=====================================================================
\section{The reduced Keller identity for the family
\texorpdfstring{$(1,-1,-m)$}{(1,-1,-m)}}
\label{sec:redw}
%=====================================================================

The reduction used for the Alp\"oge--Mathew system in Section
\ref{sec:equivariance} holds uniformly in the weight family
$(1,-1,-m)$. It is proved as an identity of differential polynomials
(sympy applied to undetermined functions --- not a spot check) in
\texttt{scripts/reduction\_113.py}, module \texttt{jcqft/reduction\_w.py}.

\begin{proposition}[Reduced Keller identity; verified for $m\le 4$,
uniform proof]\label{prop:redw}
Fix $m \ge 1$ and the source action
$\lambda.(x,y,z) = (\lambda x, \lambda^{-1}y, \lambda^{-m}z)$, with free
invariant ring $\C[w,v]$, $w = xy$, $v = x^m z$. Invertibility of
$DF(0)$ forces the component weights of an equivariant $F$ to be a
permutation of $(-m,-1,1)$; fixing the order,
\[
F \;=\; \Bigl(\frac{P(w,v)}{x^{m}},\ \frac{Q(w,v)}{x},\ x\,R(w,v)\Bigr).
\]
Then:
\begin{enumerate}
\item $F$ is a polynomial map iff every monomial $w^jv^k$ of $P$
 satisfies $j + mk \ge m$ and every monomial of $Q$ satisfies
 $j + mk \ge 1$ ($R$ is unconstrained);
\item $\det DF$ is a function of $(w,v)$ alone:
 \[
 \det DF \;=\; \det M \;=\;
 -m\,P\,J_2(Q,R) + Q\,J_2(P,R) + R\,J_2(P,Q),
 \]
 where $M$ is the $3\times3$ matrix with rows
 $(d_iS_i,\ \partial_wS_i,\ \partial_vS_i)$,
 $(S_1,S_2,S_3) = (P,Q,R)$, $d = (-m,-1,1)$;
\item \emph{compact form:} $R^m \det M = J_2(P R^m,\, Q R)$
 identically, so
 \[
 \det DF \equiv \kappa
 \quad\Longleftrightarrow\quad
 J_2(P R^m,\, Q R) = \kappa\, R^m .
 \]
\end{enumerate}
The 3D Keller problem for this equivariance class is therefore a 2D
problem, for every $m$.
\end{proposition}

\begin{proof}[Proof sketch]
In the coordinates $(x,w,v)$ one has
$\det \partial(x,w,v)/\partial(x,y,z) = x^{m+1}$, and the equivariant
form gives $\det DF = x^{\sum d_i + m + 1}\det M = \det M$ since
$\sum d_i = -m$. The script proves the identity (and the compact form)
with $P,Q,R$ undetermined functions for $m = 1,2,3,4$; the chain-rule
computation is uniform in $m$. Consistency at $m=2$: the machinery
reproduces the known reduction on the Alp\"oge--Mathew map, with
residual $0$ at $\kappa = -2$. The polynomiality box is verified in both
directions on exhaustive monomial boxes at $m=3$.
\end{proof}

Two gauge structures accompany the reduction and are used throughout
Section \ref{sec:113}. First, $DF(0) = \antidiag(p_1(0), q_0'(0),
r_0(0))$, where $p_1$ is the $v$-coefficient of $P$, etc.; target
scalings normalize $p_1(0) = q_0'(0) = r_0(0) = 1$, hence (at $m=3$)
$\kappa = \det DF(0) = -1$. Second, the source shears
$z \mapsto z + y^m h(xy)$ act on the invariants as
$v \mapsto v + w^m\,h(w)$ --- the box-preserving part of the full
\emph{$v$-shear group} $v \mapsto v + h(w)$, $h \in \C[w]$ arbitrary,
under which $\det M$ is exactly invariant (asserted). The full $v$-shear
group is an analysis tool: one solves modulo shears first and imposes the
polynomiality box last.

%=====================================================================
\section{The weight system \texorpdfstring{$(2,-1,-3)$}{(2,-1,-3)}:
\texorpdfstring{$A_1$}{A1}-cone reduction and degree-1
classification}\label{sec:213}
%=====================================================================

All claims in this section are proved by assertion in
\texttt{scripts/search\_213.py} and its module
\texttt{jcqft/reduction\_213.py}; the repository document
\texttt{docs/SEARCH\_213.md} gives the full account.

\subsection{Why this weight system}

For $\lambda.(x,y,z) = (\lambda^2x, \lambda^{-1}y, \lambda^{-3}z)$ the
orbit space has \emph{two} orbifold strata instead of the single one of
the Alp\"oge--Mathew system: the target $x$-axis has stabilizer $\Z_2$
(the same $2{:}1$ mechanism that powers the counterexample) and the
target $z$-axis has stabilizer $\Z_3$ (a free orbit mapping into it
would be $3{:}1$ --- cube-root escape and a $\Z_3$ monodromy class, a
candidate for a different global-anomaly class). Unlike the
$(1,-1,-m)$ family the invariant ring is not free, so the reduction of
Proposition \ref{prop:redw} does not apply and was rebuilt.

\subsection{Invariant theory and normal form (exact)}

Monomials $x^iy^jz^k$ are invariant iff $2i - j - 3k = 0$; the invariant
ring is the $A_1$ quadric cone
\[
R \;=\; \C[u_1,u_2,u_3]/(u_2^2 - u_1u_3),
\qquad
u_1 = xy^2,\quad u_2 = x^2yz,\quad u_3 = x^3z^2,
\]
verified through a unique normal form for invariant monomials
(exhaustively for degree $\le 12$). The spaces of weight-$d$
polynomials are $R$-modules with verified syzygy-free splittings, e.g.\
$M_{-1} = yR \oplus xz\,\C[u_3]$ and $M_{-3} = zR \oplus y^3\C[u_1]$.
Invertibility of $DF(0)$ forces the component weights to be a
permutation of $(2,-1,-3)$, and every candidate has the normal form
\[
F \;=\; \bigl(\,x\,A(u),\ \ y\,B(u) + xz\,C(u_3),\ \
 z\,D(u) + y^3E(u_1)\,\bigr),
\qquad
DF(0) = \mathrm{diag}\bigl(A(0), B(0), D(0)\bigr),
\]
with no polynomiality side conditions (every polynomial choice of
$(A,B,C,D,E)$ assembles to a polynomial map).

\subsection{Chart-based reduced Keller identity (exact)}

The $\C^*$-action is free on $\{y\ne 0\}$, and on $\{x\ne0,\,y\ne0\}$
the slice parametrization
\[
\Phi(t,p,q) \;=\; \Bigl(\frac{p}{t^2},\ t,\ \frac{q\,t^3}{p^2}\Bigr),
\qquad t = y,\ p = u_1,\ q = u_2,
\]
pulls the invariants back to $(p, q, q^2/p)$ --- the cone in its chart
$u_3 = u_2^2/u_1$. Writing $\tilde a = A(p,q,q^2/p)$,
$\tilde b = B(\cdot) + \tfrac{q}{p}C(q^2/p)$,
$\tilde e = \tfrac{q}{p^2}D(\cdot) + E(p)$, one gets
$F\circ\Phi = (p\tilde a/t^2,\ t\tilde b,\ t^3\tilde e)$ and, by the
chain rule,
\begin{equation}\label{eq:213identity}
\det DF \;=\; \Delta(p,q) \;=\;
p^2\Bigl(2\,p\tilde a\,J_2(\tilde b,\tilde e)
 + \tilde b\,J_2(p\tilde a,\tilde e)
 - 3\,\tilde e\,J_2(p\tilde a,\tilde b)\Bigr),
\end{equation}
with $J_2$ the Jacobian in $(p,q)$: a rational function of the two
invariant chart coordinates alone, with denominator a pure power of
$p$. The identity is proved with undetermined functions and re-verified
against the raw 3D determinant and against an intrinsic normal-form
determinant on the cone. Since $\Phi$ parametrizes a dense open set,
$\det DF \equiv \kappa$ iff
$\operatorname{numer}(\Delta-\kappa) \equiv 0$ in $\C[p,q]$ --- a finite,
\emph{trilinear} polynomial system (degree $\le 1$ in each of the blocks
$A \mid B,C \mid D,E$). Gauge: target scalings set $A(0)=B(0)=D(0)=1$,
$\kappa = 1$; the constant $xz$-shear sets $C(0)=0$ (box-preservingly);
$E(0)$ cannot be gauged away and is kept as an unknown.

The two stabilizer-jump mechanisms become pointwise conditions: a free
orbit over $(p_0,q_0)$, $p_0\ne0$, maps $2{:}1$ onto the target
$x$-axis iff $\tilde b = \tilde e = 0 \ne \tilde a$ there, and $3{:}1$
onto the target $z$-axis iff $\tilde a = \tilde b = 0 \ne \tilde e$;
analogous conditions cover the $y=0$ stratum, and the $x=0$ stratum
supports no mechanism. Any Keller solution with such a witness would
automatically be a counterexample.

\subsection{The degree-1 box: complete classification}

\begin{theorem}[Degree-1 classification; exact, box-complete]
\label{thm:213class}
In the box $A,B,D$ of degree $\le 1$ in $(u_1,u_2,u_3)$, $C = c_3u_3$,
$E = e_0 + e_1u_1$ (canonical, syzygy-free), with the gauge above ---
12 unknowns, 20 trilinear equations --- the Keller variety is
2-dimensional and decomposes exactly (not just up to closure) as
$V(I) = V_1 \cup V_2$:
\begin{enumerate}
\item $V_1$ (elementary $z$-shears):
 $F = (x,\ y,\ z + y^3(e_0 + e_1u_1))$, i.e.\
 $z \mapsto z + y^3h(xy^2)$;
\item $V_2$ ($w$-shears): with $w = xz - \nu y$ and parameters
 $(\lambda,\nu)$,
 \[
 F = \bigl(x,\ y + \lambda x w^3,\ z + \lambda\nu w^3\bigr),
 \qquad w\circ F = w,
 \]
 an elementary shear in the coordinates $(x,w,z)$.
\end{enumerate}
Both families are tame automorphisms with explicit polynomial inverses
(verified by composition) and generic fiber cardinality $1$. Moreover
all six orbifold-mechanism systems ($2{:}1$ and $3{:}1$, at the
torus-normalized witness points and on the $y=0$ stratum) are provably
empty (Gr\"obner basis $=[1]$).
\end{theorem}

The completeness proof is exact: six relations have powers inside the
ideal $I$ (Gr\"obner reduction certificates), so they vanish on the
whole variety, and the ideal element $c_3e_1^3 \in I$ splits the variety
into $\{c_3\ne0\}$ (pinned to $V_2$), $\{e_1\ne0\}$ (pinned to $V_1$)
and $\{c_3 = e_1 = 0\}$ (pinned to $V_1$ again). The residual 2-torus
acts freely on witness points and preserves box and gauge, so the
pointwise mechanism checks exhaust all witnesses.

In the enlarged \emph{CE2 box} ($C \in \mathrm{span}\{u_3,u_3^2\}$,
$E \in \mathrm{span}\{1,u_1,u_1^2\}$; 14 unknowns, 34 equations) the six
mechanism emptiness queries are again all empty (exact, over $\Q$;
longest single Gr\"obner basis $\approx 5$ min). A full variety
decomposition of CE2 was not attempted, so in that box only the
stabilizer-jump mechanisms are excluded --- a Keller map could in
principle identify two distinct free orbits; ruling that out needs the
full decomposition, done only for the degree-1 box. (For
Alp\"oge--Mathew the realized mechanism \emph{is} the stabilizer jump,
which motivates prioritizing it.)

\subsection{The degree-2 box: an honest wall}\label{sec:213wall}

The full degree-2 box (29 unknowns, 57 Keller equations plus witness
conditions) is \textbf{unresolved, not ruled out}, after a serious
tooling campaign whose obstruction is now precisely mapped
(\texttt{scripts/deg2\_213\_elim.py}):
\begin{enumerate}
\item sympy's Buchberger: no query finishes ($>25$ min each, also mod
 $p$).
\item msolve F4 \cite{msolve} (under a 16 GB address-space cap): all six
 queries exceed the cap over $\Q$ --- and, decisively, also mod three
 independent $\sim$30-bit primes, even after the exact reductions
 below; the memory blow-up is intrinsic to F4 on this system, not a
 coefficient-growth artifact.
\item Exact structural reductions (all asserted): the 57 Keller
 equations are triangular in the $A$-block with integer pivots in
 $\{2,\dots,7\}$, so the $A$-block eliminates by a global polynomial
 substitution valid over any field of characteristic $0$ or $p > 7$
 (49 equations, 21 unknowns); the Rabinowitsch saturation variable is
 unnecessary because $\Delta = \tilde aX + \tilde bY + \tilde eZ$
 identically; witness equations eliminate 1--2 more unknowns. Net:
 19--20 unknowns per query --- still past the 16 GB cap for F4, even
 mod $p$.
\item Singular's truncated \texttt{std()} ladder \cite{Singular}
 (degree bound; one-sided: finding $1$ below the bound is an exact
 certificate): mod $p$, degree bounds 4 and 5 complete in under a
 minute with no certificate; bounds 6 and 7 exceed 20-minute budgets.
 Hence the unit certificate, if the ideal is the unit ideal, has
 degree $\ge 6$.
\end{enumerate}
Plausible routes forward: much larger memory, further stratification,
or numerical algebraic geometry (homotopy continuation) for discovery
with exact certification of any find; see Section \ref{sec:outlook}.

\subsection{Verdict}

The weight system $(2,-1,-3)$ \emph{does} support nonlinear Keller maps
--- two 2-parameter families --- but every Keller map in the degree-1 box
is a tame automorphism, and both orbifold mechanisms are provably empty
in the degree-1 and CE2 boxes. The $\Z_3$ stratum made this the most
promising place to look for a new global-anomaly class; the answer in
the searched boxes is a clean no-go. The contrast with $(1,-1,-2)$,
where the Alp\"oge--Mathew defect exists already at low degree, suggests
the defect needs the free invariant ring of the $(1,-1,-m)$ family, or
larger boxes; the $A_1$ cone singularity appears to rigidify the
low-degree Keller moduli. (Interpretation, clearly flagged as such.)

%=====================================================================
\section{The weight system \texorpdfstring{$(1,-1,-3)$}{(1,-1,-3)}:
complete classification of the \texorpdfstring{$v$}{v}-linear
class}\label{sec:113}
%=====================================================================

This is the $m=3$ case of the uniqueness theorem
(Section \ref{sec:11m}). Every identity and every emptiness claim
is proved by assertion in \texttt{scripts/search\_113.py} (default run
$\sim$20 min; \texttt{--full} $\sim$1 h adds larger Gr\"obner boxes);
the reduction machinery is Proposition \ref{prop:redw} at $m=3$.

\subsection{Setup: why \texorpdfstring{$(1,-1,-3)$}{(1,-1,-3)}, and the
shape of the problem}

For $\lambda.(x,y,z) = (\lambda x, \lambda^{-1}y, \lambda^{-3}z)$ the
invariant ring is free --- $\C[w,v]$ with $w = xy$, $v = x^3z$ ---
exactly as in the Alp\"oge--Mathew case $m=2$. The target coordinates
carry weights $(-3,-1,1)$, so the target $a$-axis has stabilizer
$\Z_3$: a free orbit mapping into it is $3{:}1$ (cube-root escape,
$\Z_3$ monodromy). Moreover \emph{no $2{:}1$ mechanism exists at all}
in this weight system: every other coordinate axis and every coordinate
pair has trivial stabilizer, so the only stabilizer-jump
non-injectivity available is the $3{:}1$ one. Since the
Alp\"oge--Mathew counterexample is exactly a stabilizer-jump ($2{:}1$)
map in the $m=2$ member of the same family, $(1,-1,-3)$ is the natural
place to ask whether the mechanism \emph{transplants}, with the residual
$\Z_2$ upgraded to $\Z_3$.

By Proposition \ref{prop:redw}, every candidate is
$F = (P/x^3,\ Q/x,\ xR)$ with polynomiality box $j + 3k \ge 3$ for the
monomials $w^jv^k$ of $P$ and $j + 3k \ge 1$ for those of $Q$; the gauge
$p_1(0) = q_0'(0) = r_0(0) = 1$ fixes $\kappa = -1$.

The class searched is the \emph{$v$-linear class}: $P, Q, R$ of degree
$\le 1$ in $v$ and \emph{arbitrary} degree in $w$,
\[
P = p_0(w) + p_1(w)\,v,\qquad
Q = q_0(w) + q_1(w)\,v,\qquad
R = r_0(w) + r_1(w)\,v,
\]
with box/invertibility constraints $p_0 \in w^3\C[w]$,
$q_0 \in w\,\C[w]$, $p_1(0) = q_0'(0) = r_0(0) = 1$. This is the exact
analogue of the class containing the Alp\"oge--Mathew map at $m=2$
(whose $P,Q,R$ are all $v$-linear). Since $\det M + 1$ has $v$-degree
$2$, the Keller condition is three ODE-type equations in $\C[w]$,
\[
E_2 = 0,\qquad E_1 = 0,\qquad E_0 = \kappa = -1,
\]
with (all asserted from $\det M$, not copied from $m=2$)
\[
E_2 = 2\bigl[\,q_1(p_1r_1)' - 2(p_1r_1)q_1'\,\bigr],
\qquad
E_0 = -3p_0(q_0'r_1 - q_1r_0') + q_0(p_0'r_1 - p_1r_0')
 + r_0(p_0'q_1 - p_1q_0'),
\]
and $E_1$ the mixed analogue,
\[
E_1 = -3p_0[Q,R]_1 - 3p_1[Q,R]_0 + q_0[P,R]_1 + q_1[P,R]_0
 + r_0[P,Q]_1 + r_1[P,Q]_0,
\]
where $[A,B]_i$ denotes the $v^i$-part of $J_2(A,B)$. Unlike the
degree-boxed searches of Section \ref{sec:213}, the $w$-direction here
is \emph{unbounded}: the classification below is a theorem for the whole
class, not for a box --- with one precisely-mapped exception (Section
\ref{sec:gap}).

\subsection{The Wronskian stratification}

\begin{lemma}[Integration certificate; asserted]\label{lem:wronskian}
$\bigl(p_1r_1/q_1^2\bigr)' = E_2/(2q_1^3)$. Hence on
$\{E_2 = 0,\ q_1p_1r_1 \ne 0\}$ the rational function $p_1r_1/q_1^2$ has
vanishing derivative, so
\[
p_1 r_1 \;=\; c\,q_1^2, \qquad c \in \C^\times,
\]
exactly; and $q_1 = 0$ or $r_1 = 0$ satisfy $E_2$ identically.
\end{lemma}

Since $p_1(0)=1$ forces $p_1 \ne 0$, the $v$-linear Keller problem
stratifies as follows. In stratum D, the UFD solution of
$p_1r_1 = cq_1^2$ is $p_1 = as^2g$, $q_1 = stg$, $r_1 = bt^2g$ with
$g = \gcd(p_1,r_1)$, $\gcd(s,t)=1$, $ab = c$; $E_0$ is linear in
$(p_1,q_1,r_1)$ with no derivatives of them, so $g \mid E_0 = \kappa$
and $g$ is a unit (absorbed).

\begin{center}
\small
\begin{tabular}{lll}
\hline
stratum & definition & verdict at $m=3$ \\
\hline
A & $q_1 = r_1 = 0$ & tame family (all $w$-degrees) \\
B & $r_1 = 0,\ q_1 \ne 0$ & tame family (all $w$-degrees) \\
C & $q_1 = 0,\ r_1 \ne 0$ & empty (all $w$-degrees) \\
D0 & $s,t$ const ($p_1,q_1,r_1$ const) & empty (all degrees) \\
D1 & $t$ const, $s$ nonconst (\emph{the AM-analogue}) &
 empty (all degrees) \\
D2 & $s$ const, $t$ nonconst & empty (all degrees) \\
D3 & $s,t$ nonconst, coprime, both squarefree & empty (all degrees) \\
$D_3'$ & $s$ or $t$ non-squarefree ($\deg p_1 \ge 4$ or $\deg r_1\ge4$) &
 empty in Gr\"obner boxes; open beyond \\
\hline
\end{tabular}
\end{center}

\emph{Cross-check at $m=2$.} The same $E_2$-integration at $m=2$ gives
$p_1^2r_1 = c\,q_1^3$, and the Alp\"oge--Mathew data satisfies it with
$c = -1/27$: $p_1 = (1+w)^3 = s^3$, $q_1 = 3(1+w)^2 = 3s^2$, $r_1 = -1$
(asserted). That is, the Alp\"oge--Mathew map lives in the $m=2$
analogue of stratum D1 --- the stratum that is \emph{empty} at $m=3$.
The machinery does not spuriously kill the known counterexample; the two
weight systems genuinely differ.

\subsection{The main theorem}

\begin{theorem}[Complete classification of the $v$-linear class,
$m=3$; exact, all $w$-degrees, modulo the $D_3'$ gap of Section
\ref{sec:gap}]\label{thm:113main}
In the gauge $p_1(0) = q_0'(0) = r_0(0) = 1$, $\kappa = -1$, the
complete solution set of the Keller condition in the $v$-linear class
is the one-function-plus-one-parameter family
\[
\boxed{\ P = p_0 + v,\qquad Q = w + b_0P,\qquad R = 1,\qquad
p_0 \in w^3\,\C[w],\ \ b_0 \in \C .\ }
\]
Every member assembles to
$F = \bigl(p_0(xy)/x^3 + z,\ \ y + b_0x^2F_1,\ \ x\bigr)$: an
elementary $z$-shear followed by the target shear
$b \mapsto b + b_0\,ac^2$ --- a \emph{tame automorphism}, with explicit
polynomial inverse verified by composition and generic fiber
cardinality $1$. Strata C, D0, D1, D2 and D3 (with $s,t$ squarefree)
are empty for every $w$-degree.
\end{theorem}

\begin{proof}[Proof sketch]
We go stratum by stratum. Each step combines an identity chain
asserted in \texttt{scripts/search\_113.py} with a one-line
divisibility argument in $\C[w]$ (a nonzero constant has no
nonconstant divisors).

\emph{A.} $E_1 \equiv 0$ and $E_0 = -p_1(q_0r_0)' = \kappa$, so
$p_1 \mid \kappa$, giving $p_1 = 1$ (gauge); then $q_0r_0 = w$ with
$q_0(0) = 0$, $r_0(0) = 1$ forces $r_0 = 1$, $q_0 = w$; $p_0 \in
w^3\C[w]$ stays free.

\emph{B.} $E_1$ integrates ($(p_1r_0^2/q_1)' = r_0E_1/q_1^2$, asserted)
to $p_1r_0^2 = c\,q_1$; then
$E_0 = p_1\bigl[(p_0r_0^3)/c - q_0r_0\bigr]'$, forcing $p_1 = 1$,
$r_0 \mid w$ with $r_0(0)=1$ so $r_0 = 1$, and $q_0 = b_0p_0 + w$,
$q_1 = b_0 = 1/c$. Together with A (the case $b_0 = 0$) this gives the
boxed family.

\emph{C.} $(p_1/(q_0^4r_1))' = E_1/(q_0^5r_1^2)$ (asserted), so
$E_1 = 0$ forces $p_1 = c\,q_0^4r_1$ and hence $p_1(0) = 0$,
contradicting $p_1(0) = 1$.

\emph{D0} ($p_1,q_1,r_1$ nonzero constants $\alpha,\beta,\gamma$):
$E_1$ integrates to a linear relation; substituting into $E_0$ gives an
exact derivative $[-(2\alpha\gamma/\beta)V^2 - \alpha\delta V]' =
\kappa$ with $V = q_0 - (\beta/\alpha)p_0$, $V(0)=0$, $V'(0)\ne0$. Then
$V \mid \kappa w$ forces $V = \varepsilon w$, whose quadratic term
forces $2\alpha\gamma/\beta = 0$: impossible.

\emph{D1} ($p_1 = As^2$, $q_1 = Bs$, $r_1 = C$, $s$ nonconstant --- the
Alp\"oge--Mathew slot): the $v$-shear $h = -r_0/C$ sets $r_0 \equiv 0$;
then $E_1$ is a linear ODE for $p_0$ with particular solution
$p_0^* = (2A/B)sq_0$ and homogeneous solutions satisfying
$Y^2 = e\,s^3$ (nonzero only if $s$ is a constant times a square). If
$s$ is not: $E_0(p_0^*) = (2AC/B)\,q_0(s'q_0 - 2sq_0') = \kappa$ forces
$q_0 = n$ constant and $s$ linear. If $s = cd^2$: $E_0$ factors with a
factor $d \mid \kappa$, which kills it. The surviving sheared family
$(s, q_0, p_0, r_0) = (s_0 + s_1w,\ n,\ (2A/B)ns,\ 0)$ has
$\kappa = (2AC/B)n^2s_1 \ne 0$, so $n \ne 0$; but \emph{un-shearing
into the box} ($p_0(0) = q_0(0) = 0$) forces $n = 0$. Empty.

\emph{D2} ($p_1 = A$ constant): the shear $h = -p_0/A$ sets
$p_0 \equiv 0$; $E_0 = -A(q_0r_0)' = \kappa$ makes $q_0r_0$ exactly
linear, so one factor is constant. If $q_0 = \gamma$: $E_1$ forces $t$
linear and $r_0 = (2C\gamma/B)t$, but un-shearing needs
$\val_w(p_0 = -Ah) \ge 3$, hence $h(0) = 0$, hence $q_0(0) = \gamma \ne
0$: contradiction with $q_0(0)=0$. If $r_0 = \delta$: $E_1$ gives
$t \mid \delta^2t'$, so $t' = 0$: contradiction.

\emph{D3} ($s,t$ nonconstant, coprime, squarefree): pass to the
shear-invariant combinations $G_1 = tp_0 - asq_0$,
$G_2 = btq_0 - sr_0$. Asserted identities express $E_1$ and a
combination of $E_0, E_1$ through $G_1, G_2$ and
$G_3 = btG_1 + asG_2$; $G_1 = 0$ or $G_2 = 0$ contradict
$\kappa \ne 0$. Squarefree divisibility gives $s \mid G_1$,
$t \mid G_2$; writing $Z = a\bar g_2 - b\bar g_1$ and the Wronskian
$W = s't - st'$, one finds $E_1 = st(WZ - 2stZ')$ and
$(Z^2t/s)' \propto Z(2stZ' - WZ)$, forcing $Z^2t = es$ and hence
$Z = 0$. The $\kappa$-equation then reads
$g(2stg' - Wg) = -b\kappa/(2a)$, so $g$ is a constant $c$ and
$\kappa = (2a/b)c^2W$: Keller solutions exist iff the Wronskian $W$ is
a nonzero constant, and they form the single shear-orbit
$p_0 = as(q_0 + c/b)/t$, $r_0 = t(bq_0 - c)/s$. But $tp_0 =
as(q_0 + c/b)$ evaluated at $w=0$ gives $p_0(0)t(0) = a\,s(0)\,c/b \ne
0$, while the box forces $p_0(0) = 0$: empty for every degree.
\end{proof}

\begin{remark}
The $v$-shear moves (arbitrary $h \in \C[w]$) are \emph{not}
box-preserving; the box is imposed at the end through the jet
conditions at $w = 0$. This is exactly where the $m=3$ obstruction
lives (Section \ref{sec:numerology}).
\end{remark}

\subsection{The gap: \texorpdfstring{$D_3'$}{D3'}, and its closure in
boxes}\label{sec:gap}

The two steps $s \mid s'G_1 \Rightarrow s \mid G_1$ and
$t \mid t'G_2 \Rightarrow t \mid G_2$ in the D3 argument need $s,t$
squarefree. A non-squarefree $s$ or $t$ requires $\deg p_1 = 2\deg s
\ge 4$ or $\deg r_1 = 2\deg t \ge 4$. This corner is closed by exact
in-box Gr\"obner certificates (msolve \cite{msolve}, over $\Q$, under a
16 GB memory cap):
\begin{enumerate}
\item \emph{Default run:} Keller $+$ ($r_1$ has a nonzero coefficient)
 is the unit ideal in the box
 $\deg(p_0,p_1,q_0,q_1,r_0,r_1) \le (4,2,3,2,2,2)$ --- covering strata
 C and D wholesale ($\sim$5 s per coefficient).
\item \emph{\texttt{--full}:} the same in the medium box
 $\le (5,3,4,3,3,3)$, plus the four targeted non-squarefree
 parametrizations $s = (1+\rho w)^2$ / $t = (n_0 + n_1w)^2$ with the
 partner of degree $\le 2$ and $\deg(p_0,q_0,r_0) \le (6,5,4)$, each a
 direct emptiness query in 18--19 unknowns. Status (2026-07-24, 30 GB
 machine): the $s$-square/$t$-linear query is \textbf{empty} (exact,
 649 s); the remaining three hit msolve's 16 GB F4 memory wall --- the
 same wall as Section \ref{sec:213wall} --- and remain formally
 unresolved in the large box. They \emph{are} covered by the
 medium-box $r_1 \ne 0$ queries up to $\deg r_1 \le 3$, so only the
 $\deg = 4$ corner is genuinely open.
\end{enumerate}
Beyond the boxes, $D_3'$ is \emph{open}, though unpromising: the
squarefree D3 analysis shows the entire stratum's solution set is one
shear-orbit killed by the origin jet, and the non-squarefree variant
only tightens the divisibility constraints. Independently of the gap,
the classification is corroborated in-box by nilpotency certificates:
on the $r_1 = 0$ slice of the small box, powers of the relevant
coefficient combinations all lie in the Keller ideal, pinning the
in-box variety exactly to the gauged family of Theorem
\ref{thm:113main} (with its box truncation).

\subsection{The \texorpdfstring{$3{:}1$}{3:1} mechanism is empty; no
\texorpdfstring{$2{:}1$}{2:1} exists}

\begin{corollary}[Verified]\label{cor:31}
A $3{:}1$ stabilizer-jump witness needs a point $(w_0,v_0)$ with
$Q = R = 0 \ne P$ on a free orbit $\{x \ne 0\}$ (the $x=0$ plane
supports no mechanism: $F_3 \equiv 0$ there and the induced 2D map is
linear in its orbit parameter). By Theorem \ref{thm:113main} every
$v$-linear Keller map has $R = 1$ in the gauge --- $R$ never vanishes,
so no witness exists, for any $w$-degree. Independent pointwise
Gr\"obner queries (the residual torus moves any witness to $(1,1)$,
$(1,0)$ or $(0,1)$; $(0,0)$ is excluded by $R(0,0)=1$) confirm
emptiness in the small box. By weight arithmetic there is no $2{:}1$
mechanism anywhere in this weight system, so the $v$-linear class of
$(1,-1,-3)$ admits no orbifold-type counterexample at all.
\end{corollary}

The family members survive the repository's non-properness prefilter
only through its documented false-positive class of nonlinear
automorphisms (module \texttt{jcqft/prefilter.py}), and injectivity is
re-checked directly: explicit polynomial inverses by composition, and
generic fiber cardinality $1$ on samples.

\textbf{Headline: no counterexample.} Every Keller map of the
$v$-linear class is a tame automorphism.

\subsection{Why Alp\"oge--Mathew does not transplant: the exact
numerology}\label{sec:numerology}

At $m=2$ the Alp\"oge--Mathew map sits in stratum D1 of its own family:
$p_1 = s^3$, $q_1 = 3s^2$, $r_1 = -1$, $s = 1+w$ --- the same
``$t$ constant'' slot that at $m=3$ produces the sheared family
$q_0 = n$, $p_0 = (2A/B)ns$. The $m=3$ box demands $\val_w p_0 \ge 3$
(from $j + 3k \ge 3$), while the family jets force $p_0(0) \ne 0$
unless $n = 0$, i.e.\ $\kappa = 0$. At $m=2$ the box demands only
$\val_w p_0 \ge 2$, \emph{and} the corresponding equations carry a
different derivative-weighting ($E_0^{m=2} = q_0p_0' - 2p_0q_0' +
\dots$), which is exactly what lets
$P_{\rm AM} = (1+w)\bigl(v(1+w)^2 + w^2(3w+4)\bigr)$ thread the needle.
The obstruction at $m=3$ is \emph{numerological and exact} --- not a
failure to search far enough.

\subsection{Honest limitations}

\begin{enumerate}
\item \emph{$v$-degree.} The class searched is $v$-linear (the
 AM-analogue class). Ans\"atze of $v$-degree $\ge 2$ are untouched;
 for them the Keller condition has higher $v$-degree and the Wronskian
 stratification does not directly apply. (At $m=2$ the known
 counterexample is $v$-linear; a heuristic, not an argument, for
 prioritizing the class.)
\item \emph{The $D_3'$ gap} (Section \ref{sec:gap}): closed only in
 Gr\"obner boxes; open beyond, with the $\deg = 4$ corner the
 genuinely open part.
\item \emph{Gauge bookkeeping.} The classification is stated in the
 gauge $p_1(0)=q_0'(0)=r_0(0)=1$, $\kappa=-1$; un-gauged solutions are
 recovered by target/torus scalings, which change neither tameness,
 properness nor fiber counts.
\end{enumerate}

%=====================================================================
\section{Uniqueness for every $m\ge 3$}\label{sec:11m}
%=====================================================================

The $m=3$ classification of Section \ref{sec:113} extends uniformly.
Every identity below is asserted in \texttt{scripts/search\_11m.py}
(78 checks, $\sim$1\,min default; \texttt{--full} $\sim$72\,min), with
$m$ a sympy \texttt{Symbol} wherever the statement is uniform and with
$k$-symbolic parity branches ($m=2k$ / $2k{+}1$) otherwise; spot-checks
at $m=2,3,4,5,7$. Full write-up: \texttt{docs/SEARCH\_11M.md}.

\begin{theorem}[Uniqueness in the $v$-linear class; verified with gaps]
\label{thm:11m}
For every integer $m\ge 3$, every Keller map in the $v$-linear class of
$(1,-1,-m)$ is, up to gauge, the tame automorphism
\[
P = p_0 + v,\qquad Q = w + b_0 P,\qquad R = 1,\qquad
p_0\in w^m\,\C[w],\quad b_0\in\C.
\]
Hence that class contains no counterexample and no $m{:}1$ orbifold
covering for any $m\ge 3$, and the Alp\"oge--Mathew map ($m=2$) is the
unique member of its equivariant family within the $v$-linear class.
The same $D_3'$ non-squarefree gap as at $m=3$ remains, box-closed for
$m\le 5$.
\end{theorem}

\emph{The general integrated constraint.} The $v^2$-coefficient of
$\det M$ (with $m$ symbolic) is
$E_2 = 2p_1'q_1r_1 - (m{+}1)p_1q_1'r_1 + (m{-}1)p_1q_1r_1'$,
i.e.\ log-derivative exponents $(2,-(m{+}1),m{-}1)$. Integration gives
\[
p_1^2\, r_1^{m-1} = c\, q_1^{m+1}\qquad (c\in\C^\times),
\]
via the asserted certificate
$(p_1^2 r_1^{m-1}/q_1^{m+1})' = p_1 r_1^{m-2} E_2 / q_1^{m+2}$.
Anchors: $m=2$ recovers $p_1^2 r_1 = c q_1^3$ (AM at $c=-1/27$);
$m=3$ recovers the square of Lemma \ref{lem:wronskian}.

\emph{Stratum verdicts.} Strata A/B give the tame family and C, D0 are
empty, all with $m$ symbolic. For D1--D3 the bookkeeping branches on
parity of $m$; for $m\ge 4$ the $\kappa$-equations carry nonconstant
killing factors --- $u^{m-2}$ (D1 even), $u^{(m-3)/2}$ and $d^{m-2}$
(D1 odd), $s^{k-1}$ (D3 odd), $s^{m-2}$ (D3 even) --- forcing a
nonconstant polynomial to divide the constant $\kappa$. Thus for
$m\ge 4$ the D-strata are empty \emph{outright, with no box}. Only
$m=3$ needs the box-jet arguments of Section \ref{sec:113}; it is the
boundary case of the uniform mechanism.

\emph{The exact $m=2$ degeneracy.} Alp\"oge--Mathew sits in D1-even,
where two degeneracies coincide at $m=2$ and only there: (i) the
killing factor $u^{m-2}$ is trivial; (ii) the even-only homogeneous
$E_1$-direction $Y=\tilde y\, u^m$ exists, and AM needs
$\tilde y=-1\neq 0$. Asserted end-to-end: AM's sheared data
$(p_0^{\mathrm{sh}}=(w{+}1)(w{+}2),\ q_0^{\mathrm{sh}}=6{+}4w)$
plugs into the $\kappa$-formula and yields $E_0=-2=\kappa_{\mathrm{AM}}$
exactly. The machinery does not spuriously kill the known
counterexample.

\emph{Orbifold.} For every $m\ge 2$ (composite included) the only
stabilizer-jump mechanism is $m{:}1$ on the weight $-m$ axis; the
classification gives $R\equiv 1$, so it is empty. Pointwise Gr\"obner
witnesses confirm in-box at $m=4,5$.

\emph{Honest residue.} $D_3'$ beyond boxes (thresholds grow with $m$:
$\deg\ge m{+}1$ odd / $2(m{+}1)$ even) remains open. One targeted
$m=5$ gap query hit the 16 GB F4 wall; its sibling closed EMPTY
exactly, and all eight $m=4/5$ medium-box queries closed EMPTY.

%=====================================================================
\section{Implications for AQFT and pAQFT}\label{sec:qftread}
%=====================================================================

This section is the interpretive core of the paper. The underlying
facts are the exact statements of Sections
\ref{sec:obs}--\ref{sec:11m}; the contact with AQFT/pAQFT is
\emph{structural calibration}, not a theorem about $D\ge 1$ nets. We
separate what can be said from what cannot, following the ledger of
\texttt{docs/QFT\_IMPLICATIONS.md}.

\subsection{The defect and its name}

In QFT language the Alp\"oge--Mathew phenomenon is a
\emph{non-properness defect}: a global consistency condition ---
properness of the classical field map, equivalently local constancy of
the preimage count $N(J)$ --- that is strictly stronger than every
local criterion (the nonvanishing, indeed constant, Jacobian) and that
standard field-redefinition arguments implicitly assume. It is not an
anomaly in the strict cohomological sense (no symmetry is broken,
nothing is loop-generated), but it shares the defining pattern of
global anomalies: a manipulation valid by every local criterion fails
through the global structure of configuration space. Its closest
relative is the modern notion of anomalies in parameter space: dialing
the external source around a closed cycle permutes the solution
branches (the $S_3$ monodromy --- numerical evidence,
\texttt{scripts/monodromy.py}), mediated entirely through infinity in
field space. The exact observable-algebra statement of Section
\ref{sec:obs} is the algebraic face of the same defect: $F^*$ is a
monomorphism but not an automorphism of the observable algebra; the
two missing module generators $x,x^2$ are the vacuum separators; and
their separator coefficients carry the escape polynomial $p$ verbatim.

A second, related lesson is about the \emph{kind} of non-perturbative
effect at work. The standard catalogue centres on finite-action saddle
points (instantons). Here the extra sectors are not saddles at finite
field values: they live on the boundary at infinity of field space and
communicate with the finite region only across the non-properness
hypersurface $\{p=0\}$ --- transitions ``through infinity,'' realized
exactly by the escape curve of Theorem~\ref{thm:nonintegral}. Whether
an analogous boundary-of-field-space mechanism can survive in a genuine
functional-integral setting is open; the example establishes that at
the level of classical solution manifolds the mechanism exists and is
compatible with completely trivial loop structure.

\subsection{Haag--Kastler AQFT: algebras over field coordinates}
\label{sec:aqft}

Haag--Kastler AQFT \cite{HK64} takes the \emph{algebras of observables},
not field coordinates, as the primary data; fields are regarded as
interchangeable ``coordinates'' on the theory (Borchers classes). The
counterexample turns this philosophical preference into a
theorem-sized fact in the only setting we control completely. The
pullback
\[
F^*:\ \C[a,b,c]\longrightarrow\C[x,y,z],\qquad
F^*(\mathcal{O})=\mathcal{O}\circ F,
\]
is an injective algebra homomorphism that is \emph{not surjective}
(Theorems~\ref{thm:basis},~\ref{thm:nonintegral}): the polynomial
observables of the ``redefined field'' $F(\varphi)$ form a proper
subalgebra of the observables of $\varphi$, of index measured by the
degree-3 field extension with Galois closure group $S_3$. So: a
polynomial field redefinition with invertible propagator and constant
unit Jacobian can be a monomorphism, but not an automorphism, of the
observable algebra. Any framework whose objects are algebras and whose
morphisms must be checked for surjectivity is structurally protected
against the mistake this example punishes; any framework that treats
``invertible field redefinition'' as certified by the Jacobian is not.
This is, we believe, the cleanest AQFT-flavoured lesson of the
counterexample, and it is exact.

\subsection{pAQFT (Fredenhagen--Rejzner): calibrating the deferral}
\label{sec:paqft}

In perturbative algebraic QFT \cite{FR12,FR16,Rej16} interacting
observables are constructed from free ones by M{\o}ller-type maps built
as formal power series in the coupling and $\hbar$; classical
inversions of nonlinear field equations enter through the (classical)
M{\o}ller map, and invertibility holds automatically at the
formal-series level. Nothing in the counterexample exhibits an
inconsistency in this: the constructions are internally coherent as
formal deformation quantization, and their formal character is a
\emph{stated feature} of the framework. What the counterexample adds
is a calibration of the worst case of that deferral, sharper than
previously available: the formal inverse here is not merely consistent
but \emph{convergent}, analytic on a neighbourhood of the vacuum, and
satisfies every identity perturbation theory can formulate --- and
still describes only one of three solution sectors, the others being
invisible at every order because they sit at infinite field strength.
So ``upgrade formal to convergent'' is \emph{not} the missing step
between pAQFT and a non-perturbative construction; the missing input
is global (properness of the classical dynamics), and it must come
from outside perturbation theory. Relatedly, statements of
\emph{perturbative agreement} (independence of the split of the action
into free and interacting parts \cite{HW05}) are
field-redefinition-adjacent moves established as formal-series
identities; the example is a reminder of exactly which global questions
such identities do not decide.

The same point bears on field redefinitions in renormalization theory
more broadly. Arguments that treat a polynomial redefinition with
$\det DF=\mathrm{const}$ as an exact change of variables in a
functional integral or a lattice measure --- ``$\det DF=1$, so the
measure is preserved'' --- use a false step: the pushforward under a
non-injective $F$ overcounts sheets. Global injectivity is an honest
hypothesis that must be verified, not read off from the Jacobian.
Purely perturbative equivalence statements are untouched.

\subsection{Non-perturbative algebraic constructions}
\label{sec:bf}

The Buchholz--Fredenhagen $C^*$-algebraic approach \cite{BF20} builds
interacting dynamics from unitaries $S(f)$ subject to causal
factorization relations, bypassing formal series altogether. Nothing
in a 0D polynomial map bears on the correctness of that program. The
counterexample instead supplies the minimal instance of a phenomenon
any non-perturbative framework must \emph{represent} somewhere in its
data: classical dynamics whose solution count jumps (here
$1\leftrightarrow 3$) as an external source crosses an algebraic wall,
with the extra solutions entering from infinity in field space and
permuted by an $S_3$ monodromy under cycles of the source. A
well-posed exercise --- open, and we claim nothing about its outcome
--- is to formulate the 0D caricature of the $S(f)$ relations for this
$F$ and see how the wall $\{p=0\}$ and the sheet structure are
encoded.

\subsection{Rigidity of the defect: what the sibling searches say}

The $(1,-1,-3)$ grading was the closest sibling of the Alp\"oge--Mathew
theory with a $\Z_3$ orbifold stratum in place of $\Z_2$ --- the
candidate for a cube-root-escape global anomaly class. The verdict
(Sections~\ref{sec:113},~\ref{sec:11m}) is a clean no-go in the
AM-analogue class: the unit-determinant constraint is compatible with
this grading only through shear-type field redefinitions; the theory
admits no non-properness defect there and no $\Z_3$ vacuum-escape
monodromy. Combined with the $(2,-1,-3)$ no-go of Section
\ref{sec:213} and the uniqueness theorem of Section \ref{sec:11m}, the
pattern is that the Alp\"oge--Mathew defect is \emph{rigidly attached
to the $m=2$ numerology}: neither enlarging the orbifold group,
deforming the invariant geometry, nor moving to any other $m\ge 3$
reproduces it. For $m\ge 4$ the failure is an exact arithmetic
incompatibility that needs no search box at all; $m=3$ is the boundary
case. Within the $v$-linear class the defect is a delicate coincidence
of the $(1,-1,-2)$ weight arithmetic --- the killing factor $u^{m-2}$
trivializing together with the even-only homogeneous $E_1$-direction.

For the QFT reading this is essential: it means the non-properness
defect is \emph{not} a generic feature of every constant-Jacobian
theory, but a rare, exactly characterisable pathology. Frameworks that
check global invertibility of field redefinitions are not being asked
to police an open dense set of pathologies; they are being asked to
exclude a delicate coincidence that, in the equivariant family we can
control, sits only at $m=2$.

\subsection{What the model does not imply}

Plainly, and to resist overstatement:

\begin{itemize}
\item \emph{Nothing about UV renormalization.} The model is
 zero-dimensional; there are no short-distance singularities, no
 counterterms, no running. The phenomenon is invariant under everything
 renormalization theory cares about.
\item \emph{Nothing about Borel summability} of $\phi^4_2$ or
 $\phi^4_3$. Those are loop/large-order phenomena; our series is
 convergent. The two situations are disjoint.
\item \emph{Nothing about the core $D=4$ difficulties.} Suspected
 triviality of $\phi^4_4$ and the Yang--Mills existence/mass-gap
 problem are ultraviolet and measure-theoretic. The counterexample does
 not bear on them.
\item \emph{No spacetime, no net, no causality, no Hilbert space.}
 Analogies between ``extra solution sheets'' and ``sectors'' are
 structural motivation, not theorems about $D\ge 1$ AQFT.
\end{itemize}

The honest headline is narrower and still interesting:
\emph{constant-Jacobian polynomial dynamics can hide non-perturbative
sectors at infinity, and no amount of perturbative or even
locally-analytic information detects them --- but in the equivariant
family $(1,-1,-m)$ this pathology is unique to $m=2$.}

%=====================================================================
\section{Outlook}\label{sec:outlook}
%=====================================================================

Next probes, ordered by expected leverage for the QFT reading:

\begin{enumerate}
\item \emph{A 0D Buchholz--Fredenhagen caricature.} Formulate the
 $S(f)$-type relations for this $F$ and ask how the wall $\{p=0\}$ and
 the $S_3$ sheet structure are encoded --- the cleanest open exercise
 suggested by Section~\ref{sec:bf}.
\item \emph{The $v$-quadratic class of $(1,-1,-m)$.} Does anything else
 live near Alp\"oge--Mathew at $m=2$, or is uniqueness stable under
 raising the $v$-degree? Directly tests how rare the defect is.
\item \emph{Closing the $D_3'$ gap.} A structural argument removing the
 squarefree hypothesis would make Theorem~\ref{thm:11m} fully
 gap-free, strengthening the uniqueness claim that underwrites the
 ``rare pathology'' reading.
\item \emph{Four-field gradings.} Richer stabilizer patterns and
 orbifold strata; the natural place to look for a genuinely different
 global-anomaly class once the three-field family is exhausted.
\item \emph{The degree-2 $(2,-1,-3)$ box.} Still beyond direct Gr\"obner
 methods on 30\,GB hardware (Section~\ref{sec:213wall}); further
 stratification is needed before homotopy continuation is practical.
\end{enumerate}

%=====================================================================
\section{Reproducibility}\label{sec:repro}
%=====================================================================

All scripts run from the repository root as
\texttt{.venv/bin/python scripts/<name>.py}; every displayed identity
is an \texttt{assert} in the cited script. Runtimes are indicative
(single desktop machine, 30 GB RAM).

\begin{center}
\small
\begin{tabular}{lll}
\hline
script & claim(s) verified & runtime \\
\hline
\texttt{verify\_counterexample.py} &
 $\det DF \equiv -2$; triple fiber \eqref{eq:AM} & seconds \\
\texttt{branch\_locus.py} &
 eliminant cubic \eqref{eq:cubic}; wall $\{p=0\}$; exact fibers &
 seconds \\
\texttt{monodromy.py} &
 $S_3$ sheet monodromy (\emph{numerical}) & minutes \\
\texttt{missing\_observables.py} &
 Theorems \ref{thm:basis}, \ref{thm:nonintegral} & $\sim$15 s \\
\texttt{reduction\_113.py} &
 Proposition \ref{prop:redw} ($m \le 4$; boxes at $m=3$) & $\sim$4 s \\
\texttt{search\_213.py} &
 Section \ref{sec:213}: identity \eqref{eq:213identity}; Theorem
 \ref{thm:213class} & $\sim$40 s \\
\texttt{search\_213.py --full} &
 CE2-box mechanism emptiness & $\sim$8 min \\
\texttt{deg2\_213\_elim.py} &
 degree-2 box reductions and capped screens (Section
 \ref{sec:213wall}) & minutes \\
\texttt{search\_113.py} &
 Section \ref{sec:113}: Lemma \ref{lem:wronskian}, Theorem
 \ref{thm:113main}, Corollary \ref{cor:31} & $\sim$20 min \\
\texttt{search\_113.py --full} &
 medium-box and targeted $D_3'$ certificates (Section \ref{sec:gap}) &
 $\sim$1 h \\
\texttt{search\_11m.py} &
 Section \ref{sec:11m}: Theorem \ref{thm:11m} (78 assertions) &
 $\sim$1 min \\
\texttt{search\_11m.py --full} &
 medium-box / targeted $D_3'$ certificates for $m\le 5$ &
 $\sim$72 min \\
\hline
\end{tabular}
\end{center}

\subsection*{Acknowledgements}

The verification scripts and the analysis sessions behind this paper
were written and run with AI coding agents, as described in Section
\ref{sec:provenance}. The counterexample studied here is due to
Alp\"oge and Mathew \cite{AM26}.

%=====================================================================
\begin{thebibliography}{99}
%=====================================================================

\bibitem{AM26}
L.~Alp\"oge,
announcement of a counterexample to the Jacobian Conjecture in
dimension $3$ (with the question posed by A.~Mathew; the map produced
by a large language model), public announcement on X
(\texttt{@\_\_alpoge\_\_}), July 19, 2026.
% NOTE (author to update): at the time of writing no arXiv preprint is
% cited in the source repository; if/when the Alpöge–Mathew preprint
% appears, replace this item with the arXiv reference.
A Lean~4 formalization is available at
\texttt{github.com/deancureton/jacobian}.

\bibitem{Kel39}
O.-H.~Keller,
\emph{Ganze Cremona-Transformationen},
Monatsh.\ Math.\ Phys.\ \textbf{47} (1939) 299--306.

\bibitem{BCW82}
H.~Bass, E.~H.~Connell, D.~Wright,
\emph{The Jacobian conjecture: reduction of degree and formal expansion
of the inverse},
Bull.\ Amer.\ Math.\ Soc.\ (N.S.) \textbf{7} (1982) 287--330.

\bibitem{Wri87}
D.~Wright,
\emph{Formal inverse expansion and the Jacobian conjecture},
J.\ Pure Appl.\ Algebra \textbf{48} (1987) 199--219.

\bibitem{vdE00}
A.~van den Essen,
\emph{Polynomial Automorphisms and the Jacobian Conjecture},
Progress in Mathematics \textbf{190}, Birkh\"auser, Basel (2000).

\bibitem{Abd03}
A.~Abdesselam,
\emph{The Jacobian conjecture as a problem of perturbative quantum
field theory},
Ann.\ Henri Poincar\'e \textbf{4} (2003) 199--215.
arXiv:math/0208173.

\bibitem{AR95}
A.~Abdesselam, V.~Rivasseau,
\emph{Trees, forests and jungles: a botanical garden for cluster
expansions},
in \emph{Constructive Physics}, Lecture Notes in Physics \textbf{446},
Springer (1995) 7--36. arXiv:hep-th/9409094.

\bibitem{HK64}
R.~Haag, D.~Kastler,
\emph{An algebraic approach to quantum field theory},
J.\ Math.\ Phys.\ \textbf{5} (1964) 848--861.

\bibitem{FR12}
K.~Fredenhagen, K.~Rejzner,
\emph{Perturbative algebraic quantum field theory},
in \emph{Mathematical Aspects of Quantum Field Theories},
Springer (2015) 17--55. arXiv:1208.1428.

\bibitem{FR16}
K.~Fredenhagen, K.~Rejzner,
\emph{Quantum field theory on curved spacetimes: Axiomatic framework
and examples},
J.\ Math.\ Phys.\ \textbf{57} (2016) 031101. arXiv:1412.5125.

\bibitem{Rej16}
K.~Rejzner,
\emph{Perturbative Algebraic Quantum Field Theory},
Mathematical Physics Studies, Springer (2016).

\bibitem{HW05}
S.~Hollands, R.~M.~Wald,
\emph{On the renormalization group in curved spacetime},
Commun.\ Math.\ Phys.\ \textbf{237} (2003) 123--160.

\bibitem{BF20}
D.~Buchholz, K.~Fredenhagen,
\emph{A $C^*$-algebraic approach to interacting quantum field theories},
Commun.\ Math.\ Phys.\ \textbf{377} (2020) 947--969.
arXiv:1902.06062.

\bibitem{Jel93}
Z.~Jelonek,
\emph{The set of points at which a polynomial map is not proper},
Ann.\ Polon.\ Math.\ \textbf{58} (1993) 259--266.

\bibitem{dBvdE05}
M.~de Bondt, A.~van den Essen,
\emph{A reduction of the Jacobian conjecture to the symmetric case},
Proc.\ Amer.\ Math.\ Soc.\ \textbf{133} (2005) 2201--2205.

\bibitem{msolve}
J.~Berthomieu, C.~Eder, M.~Safey El Din,
\emph{msolve: A library for solving polynomial systems},
Proceedings of ISSAC 2021, ACM (2021) 51--58.

\bibitem{Singular}
W.~Decker, G.-M.~Greuel, G.~Pfister, H.~Sch\"onemann,
\textsc{Singular} --- {A} computer algebra system for polynomial
computations, \url{https://www.singular.uni-kl.de}.

\bibitem{sympy}
A.~Meurer et al.,
\emph{SymPy: symbolic computing in Python},
PeerJ Computer Science \textbf{3} (2017) e103.

\end{thebibliography}

\end{document}
